\chapter{Carpal Tunnel Syndrome}

\section*{Definition}
Carpal tunnel syndrome (CTS) is compression of the median nerve at the wrist, causing tingling, numbness, or pain in the thumb, index, middle, and thumb-side half of the ring finger. It is usually diagnosed clinically; nerve-conduction studies can support the diagnosis and assess severity when needed.

\section*{When to See a Doctor}

\begin{itemize}
 \item Persistent numbness or tingling in the median nerve distribution
 \item Nighttime symptoms that awaken you
 \item Thenar muscle weakness or atrophy
 \item Difficulty with fine motor tasks (buttoning, gripping)
\end{itemize}

\section*{How It Develops}
Increased pressure within the carpal tunnel (normally 2--10 mmHg) impairs epineural blood flow and axonal transport of the median nerve. Chronic compression leads to demyelination, perineural fibrosis, and eventual axonal degeneration. Prolonged wrist flexion or extension further elevates tunnel pressure, exacerbating symptoms.

\section*{Causes}
\begin{itemize}
 \item \textbf{idiopathic:} Most common; often bilateral with no identifiable cause
 \item \textbf{Work and activity factors:} Forceful, repetitive, or vibrating hand use may contribute, although ordinary keyboard use alone has not been established as a sole cause
 \item \textbf{Anatomic:} Small carpal tunnel, wrist fracture, tenosynovitis
 \item \textbf{Systemic:} Hypothyroidism, diabetes mellitus, rheumatoid arthritis, amyloidosis, acromegaly
 \item \textbf{Pregnancy:} Fluid retention in the third trimester (typically resolves postpartum)
 \item \textbf{Obesity:} Increased adipose tissue within the carpal canal
\end{itemize}

\section*{Related Symptoms}
\begin{itemize}
 \item Paresthesias in thumb, index, middle, and radial half of ring finger
 \item Nocturnal pain and burning sensation
 \item Thenar atrophy in advanced cases
 \item Weakness of thumb opposition and abduction
 \item Positive Tinel sign and Phalen maneuver on examination
\end{itemize}
\textbf{Important:} Bilateral CTS with atypical features should prompt evaluation for hypothyroidism, diabetes, or amyloidosis.

\section*{Medical Evaluation}
\begin{itemize}
 \item Clinical examination: Tinel sign, Phalen maneuver, thenar strength and sensation
 \item Nerve conduction studies and electromyography for confirmation and severity grading
 \item Ultrasound to measure cross-sectional area of the median nerve
 \item MRI reserved for atypical presentations or suspected structural lesions
 \item Tests for diabetes, thyroid disease, inflammatory disease, or another cause are guided by the history and examination rather than performed routinely for everyone
\end{itemize}

\section*{Treatment Options}
\begin{itemize}
 \item Activity modification and nocturnal splinting as first-line therapy
 \item Corticosteroid injection into the carpal tunnel for temporary relief
 \item Oral corticosteroids and NSAIDs do not treat the underlying nerve compression and are not routine long-term treatment
 \item Surgical carpal tunnel release (open or endoscopic) for moderate-to-severe or refractory cases
 \item Post-surgical hand therapy for grip and pinch strengthening
\end{itemize}

\section*{Self-Care and Home Management}
\begin{itemize}
 \item Neutral-position wrist splint worn at night and during aggravating activities
 \item Frequent breaks and ergonomic adjustments (keyboard height, wrist rest)
 \item Ice or cold packs to the wrist for 10--15 minutes after activity
 \item Over-the-counter NSAIDs for mild inflammatory symptoms
 \item Gentle exercises may be suggested by a hand therapist, but stop if they worsen numbness or pain
\end{itemize}

\section*{References}
\begin{thebibliography}{99}
\bibitem{carpal_tunnel_syndrome:ref1} Keith MW, Masear V, Amadio PC, et al. ``AAOS Clinical Practice Guideline: Diagnosis of Carpal Tunnel Syndrome.'' J Bone Joint Surg Am. 2009;91(1):247--248.
\bibitem{carpal_tunnel_syndrome:ref2} Padua L, Coraci D, Erra C, et al. ``Carpal Tunnel Syndrome: Clinical Features, Diagnosis, and Management.'' Lancet Neurol. 2016;15(12):1273--1284.
\bibitem{carpal_tunnel_syndrome:ref3} Middleton SD, Anakwe RE, McEachan JE. ``Carpal Tunnel Syndrome.'' BMJ. 2014;349:g6437.
\bibitem{carpal_tunnel_syndrome:ref4} Atroshi I, Gummesson C, Johnsson R, et al. ``Prevalence of Carpal Tunnel Syndrome in a General Population.'' JAMA. 1999;282(2):153--158.
\bibitem{carpal_tunnel_syndrome:ref5} American Academy of Orthopaedic Surgeons. ``Management of Carpal Tunnel Syndrome: Evidence-Based Clinical Practice Guideline.'' AAOS, 2016.
\bibitem{carpal_tunnel_syndrome:ref6} Verdugo RJ, Salinas RA, Castillo JL, et al. ``Surgical versus Non-Surgical Treatment for Carpal Tunnel Syndrome.'' Cochrane Database Syst Rev. 2008;(4):CD001552.
\end{thebibliography}
